

\begin{figure*}[t]
\centering
\begin{minipage}{0.97\textwidth}
\begin{lstlisting}[style=pythonstyle, firstnumber=1]
# =============================================================================
# Scenario: Spotify (Music Streaming Platform)
# Task: "Search for the song 'Blinding Lights' by The Weeknd and save it to
#        my existing playlist called 'Driving Vibes'."
# =============================================================================

def verify_task(initial_db_path: str, final_db_path: str) -> dict:
    """
    Verification function that compares initial and final database states
    to determine if the task was completed successfully.
    
    Returns a dictionary with task-relevant signals for LLM-as-a-Judge.
    """
    import sqlite3

    def get_conn(path):
        return sqlite3.connect(path)

    def fetchone_dict(cur, query, params=()):
        cur.execute(query, params)
        row = cur.fetchone()
        if row is None:
            return None
        cols = [c[0] for c in cur.description]
        return {cols[i]: row[i] for i in range(len(cols))}

    def fetchall_dicts(cur, query, params=()):
        cur.execute(query, params)
        rows = cur.fetchall()
        cols = [c[0] for c in cur.description]
        return [{cols[i]: row[i] for i in range(len(cols))} for row in rows]

    # Connect to both databases (read-only)
    conn_initial = get_conn(initial_db_path)
    conn_final = get_conn(final_db_path)

    try:
        cur_init = conn_initial.cursor()
        cur_final = conn_final.cursor()

        # Step 1: Identify the current user
        current_user = fetchone_dict(
            cur_final,
            "SELECT id, username FROM users WHERE username = ?",
            ("current_user",),
        )
        user_id = current_user["id"] if current_user else None
\end{lstlisting}
\end{minipage}
\caption{Synthesized verification code for Spotify environment (Part 1 / 2): Task description, helper functions, and user identification. The function takes two database paths (initial and final states) and returns task-relevant signals for verification.}
\label{fig:verification-code-part1}
\end{figure*}

\begin{figure*}[t]
\centering
\begin{minipage}{0.97\textwidth}
\begin{lstlisting}[style=pythonstyle, firstnumber=50]
        # Step 2: Locate the playlist 'Driving Vibes' in both database states
        playlist_query = """
            SELECT id, owner_user_id, name, description, created_at, updated_at
            FROM playlists WHERE owner_user_id = ? AND name = ?
        """
        playlist_initial = fetchone_dict(cur_init, playlist_query, (user_id, "Driving Vibes"))
        playlist_final = fetchone_dict(cur_final, playlist_query, (user_id, "Driving Vibes"))
        playlist_id = playlist_final["id"] if playlist_final else None

        # Step 3: Find the track 'Blinding Lights' by The Weeknd
        track_query = """
            SELECT t.id AS track_id, t.title, a.name AS artist_name
            FROM tracks t
            JOIN track_artists ta ON ta.track_id = t.id
            JOIN artists a ON a.id = ta.artist_id
            WHERE t.title = ? AND a.name = ?
        """
        track_final = fetchone_dict(cur_final, track_query, ("Blinding Lights", "The Weeknd"))
        track_id = track_final["track_id"] if track_final else None

        # Step 4: Get playlist tracks before and after agent execution
        playlist_tracks_query = """
            SELECT pt.track_id, pt.position, t.title AS track_title
            FROM playlist_tracks pt
            JOIN tracks t ON t.id = pt.track_id
            WHERE pt.playlist_id = ? ORDER BY pt.position
        """
        tracks_initial = fetchall_dicts(cur_init, playlist_tracks_query, (playlist_id,)) if playlist_id else []
        tracks_final = fetchall_dicts(cur_final, playlist_tracks_query, (playlist_id,)) if playlist_id else []

        # Step 5: Determine if 'Blinding Lights' was added to the playlist
        initial_track_ids = {t["track_id"] for t in tracks_initial}
        final_track_ids = {t["track_id"] for t in tracks_final}
        newly_added_tracks = [t for t in tracks_final if t["track_id"] not in initial_track_ids]
        blinding_lights_added = track_id in final_track_ids and track_id not in initial_track_ids

        # Step 6: Build result dictionary with task-relevant signals
        result = {
            "environment": "spotify",
            "task_description": "Search for 'Blinding Lights' by The Weeknd and save to 'Driving Vibes'",
            "playlist_exists": playlist_final is not None,
            "track_exists": track_final is not None,
            "playlist_tracks_before": len(tracks_initial),
            "playlist_tracks_after": len(tracks_final),
            "newly_added_tracks": newly_added_tracks,
            "blinding_lights_added_to_playlist": blinding_lights_added,
        }
        return result
    finally:
        conn_initial.close()
        conn_final.close()
\end{lstlisting}
\end{minipage}
\caption{Synthesized verification code for Spotify environment (Part 2 / 2): Core verification logic that queries the playlist and track tables, compares initial and final states, and returns structured signals. The key output \texttt{blinding\_lights\_added\_to\_playlist} indicates whether the target track was successfully added.}
\label{fig:verification-code-part2}
\end{figure*}
